\section*{अध्याय \ref{ch:g3:continents-and-oceans} --- ग्लोब: महाद्वीप और महासागर}

\begin{solution}{exo:g3:continents-and-oceans:1}
एशिया, अफ़्रीका, उत्तरी अमेरिका, दक्षिणी अमेरिका, अंटार्कटिका, यूरोप और
ऑस्ट्रेलिया। सबसे बड़ा एशिया है; सबसे छोटा ऑस्ट्रेलिया।
\end{solution}

\begin{solution}{exo:g3:continents-and-oceans:2}
प्रशांत, अटलांटिक, हिंद, दक्षिणी और आर्कटिक \omterm{def:g2:water-on-earth:ocean}{महासागर}। सबसे बड़ा प्रशांत
\omterm{def:g2:water-on-earth:ocean}{महासागर} है; सबसे छोटा आर्कटिक \omterm{def:g2:water-on-earth:ocean}{महासागर}।
\end{solution}

\begin{solution}{exo:g3:continents-and-oceans:3}
\omterm{def:g3:continents-and-oceans:poles}{उत्तरी ध्रुव} \omterm{def:g3:continents-and-oceans:globe}{ग्लोब} के सबसे ऊपर है; \omterm{def:g3:continents-and-oceans:poles}{दक्षिणी ध्रुव} सबसे नीचे है; \omterm{def:g3:continents-and-oceans:equator}{भूमध्य रेखा}
बीच में पूरा चक्कर लगाती है।
\end{solution}

\begin{solution}{exo:g3:continents-and-oceans:4}
यूरोप: उत्तरी। ऑस्ट्रेलिया: दक्षिणी। अंटार्कटिका: दक्षिणी। उत्तरी अमेरिका:
उत्तरी।
\end{solution}

\begin{solution}{exo:g3:continents-and-oceans:5}
\omterm{def:g3:continents-and-oceans:poles}{उत्तरी ध्रुव} के चारों ओर आर्कटिक \omterm{def:g2:water-on-earth:ocean}{महासागर} है। अंटार्कटिका के चारों ओर दक्षिणी
\omterm{def:g2:water-on-earth:ocean}{महासागर} है।
\end{solution}

\begin{solution}{exo:g3:continents-and-oceans:6}
अफ़्रीका से दक्षिणी अमेरिका: अटलांटिक \omterm{def:g2:water-on-earth:ocean}{महासागर}। ऑस्ट्रेलिया से पूर्व की ओर
दक्षिणी अमेरिका: प्रशांत \omterm{def:g2:water-on-earth:ocean}{महासागर}। अफ़्रीका से ऑस्ट्रेलिया: हिंद \omterm{def:g2:water-on-earth:ocean}{महासागर}।
\end{solution}

\begin{solution}{exo:g3:continents-and-oceans:7}
सैम \omterm{def:g3:continents-and-oceans:equator}{भूमध्य रेखा} के पास रहता है, जहाँ साल भर गरमी रहती है। बो किसी \omterm{def:g3:continents-and-oceans:poles}{ध्रुव} के
पास रहती है, जहाँ साल भर ठंड रहती है।
\end{solution}

\begin{solution}{exo:g3:continents-and-oceans:8}
$10 - 7 = 3$: उसकी उँगली 3 बार ज़मीन पर पड़ती है। यही उम्मीद की जा सकती है:
पृथ्वी के 100 में से लगभग 71 वर्ग पानी हैं, इसलिए 10 में से लगभग 7 बार उँगली
\omterm{def:g2:water-on-earth:ocean}{समुद्र} पर रुकती है। (सिर्फ़ 10 बार घुमाने पर उसे 6 या 8 भी मिल सकते थे: गिनती
एक खेल से दूसरे खेल में बदलती है।)
\end{solution}

\begin{solution}{exo:g3:continents-and-oceans:9}
\omterm{def:g3:continents-and-oceans:world-map}{विश्व मानचित्र} बनाने के लिए \omterm{def:g3:continents-and-oceans:globe}{ग्लोब} को प्रशांत \omterm{def:g2:water-on-earth:ocean}{महासागर} के बीच से जाने वाली एक
रेखा पर काटकर सपाट फैलाया गया। प्रशांत \omterm{def:g2:water-on-earth:ocean}{महासागर} का एक आधा भाग \omterm{def:g2:drawing-to-map:map}{मानचित्र} के
बाएँ किनारे पर पहुँच जाता है, दूसरा आधा दाएँ किनारे पर। असली पृथ्वी पर
प्रशांत \omterm{def:g2:water-on-earth:ocean}{महासागर} एक ही है: \omterm{def:g2:drawing-to-map:map}{मानचित्र} के दोनों किनारे \omterm{def:g3:continents-and-oceans:globe}{ग्लोब} पर एक ही रेखा हैं।
\end{solution}

\begin{solution}{exo:g3:continents-and-oceans:10}
टॉम ग़लत है। सबसे बड़े से सबसे छोटे तक \omterm{def:g3:continents-and-oceans:continent}{महाद्वीप} हैं एशिया, अफ़्रीका, उत्तरी
अमेरिका, दक्षिणी अमेरिका, अंटार्कटिका, यूरोप और ऑस्ट्रेलिया: अंटार्कटिका सिर्फ़
पाँचवाँ है, एशिया से बहुत छोटा। \omterm{def:g3:continents-and-oceans:world-map}{विश्व मानचित्र} उसे धोखा देता है क्योंकि चपटा
\omterm{def:g2:drawing-to-map:map}{मानचित्र} \omterm{def:g3:continents-and-oceans:poles}{ध्रुवों} के पास ज़मीन को बहुत खींच देता है। \omterm{def:g3:continents-and-oceans:poles}{दक्षिणी ध्रुव} के चारों ओर
फैला अंटार्कटिका \omterm{def:g2:drawing-to-map:map}{मानचित्र} के पूरे निचले भाग में एक लंबी पट्टी की तरह खिंच
जाता है। \omterm{def:g3:continents-and-oceans:globe}{ग्लोब} पर वह ज़मीन का एक गोल भूभाग है, और एशिया से बहुत छोटा।
\end{solution}
